% 15 · Physics — what membership is actually for
\section{WHAT MEMBERSHIP IS FOR}
\label{sec:physics}

\deckslides{slide 47, the isochrone and two-survey blocks.}

Every number so far has been a score against labels. This chapter closes the
loop: membership lists are not the product, they are the input to a measurement
of what the cluster \emph{is} --- its age, its distance, its chemistry --- and
those measurements are what chemical tagging has to deliver if it is to matter
for Galactic archaeology \citep{Freeman:02, BlandHawthorn:16}.

\begin{marginnote}[]
\entry{Isochrone}{A theoretical curve in the colour--magnitude diagram giving
the luminosity and colour of every mass of a single-age, single-metallicity
population. Fitting it to a cluster's stars returns age, distance and
reddening.}
\entry{Distance modulus}{$m - M = 5\log_{10} d - 5$, the difference between
apparent and absolute magnitude. It is the fitted distance in a CMD; a
0.1 mag error is a 5\% error in distance.}
\end{marginnote}

\subsection{The measurement problem}
\label{sec:measurement}

\begin{figure}[htbp]\centering
  \begin{tabular}{@{}c@{\hskip6pt}c@{}}
    \fig[0.44]{sky_ngc188.png} & \fig[0.44]{sky_ic166.png}
  \end{tabular}
  \caption{The two ends of the sample, which is why the same pipeline behaves
  differently on different clusters. Left: NGC~188 --- old ($7.6$ Gyr), well
  populated, chemically distinct from the field. Right: IC~166 --- sparse and
  contaminated, with several member rows whose \texttt{APOGEE\_ID} is blank
  (\S\ref{sec:validation}). In the first, the isochrone fit has hundreds of
  stars on a clean sequence; in the second, every contaminant moves the fit.}
  \label{fig:skies}
\end{figure}

Fitting an isochrone grid to a cluster means searching a three-dimensional
parameter space (age, metallicity, distance modulus) for the combination that
best reproduces the observed colour--magnitude diagram, with the reddening as a
fourth, partly degenerate axis. We do it with the PARSEC grids through
\texttt{ASteCA} \citep{Perren:15}, sampling the posterior with \texttt{emcee}
\citep{ForemanMackey:13}. Two properties of that fit dominate everything that
follows.

\textbf{Age, metallicity and reddening are partially degenerate.} In the
regions of the CMD our clusters occupy, an older, more metal-rich population
at larger reddening can mimic a younger, metal-poor one at smaller reddening.
The degeneracy is a property of the isochrone geometry, not of the data or the
sampler: the posterior stays broad in the degenerate direction even when every
star in the cluster is a perfect member. In the two-survey fit below, the
age posterior width is $\sigma_{\log a} \approx 0.97$ --- nearly a factor
of ten in age --- for a cluster whose distance is pinned to 0.01 mag.

\textbf{Contaminants bias the parameters in a specific direction.} Field stars
in a cluster's colour--magnitude diagram occupy the same locus as the cluster's
own stars in some mass ranges and a different one in others, and the bias they
introduce is not random. Adding field stars to a metal-poor globular cluster
lowers the inferred age; adding them to a young cluster raises it. Membership
quality is therefore a parameter-estimation problem, not a bookkeeping one.

\subsection{Distances from red clump giants}

The most robust single measurement available to us is the distance, because
red clump giants have nearly constant absolute magnitude: the median $K$-band
magnitude of the members that fall in the $J-K$ clump slice is a standard
candle with a known calibration. For the clusters where our membership is
deep enough to have red clump members, the distances agree with the
literature:

\begin{table}[htbp]
\centering
\small
\tabcolsep6pt
\begin{tabular}{@{}lccc@{}}
\toprule
cluster & $dm$ (clump) & $dm$ (literature) & residual \\
\midrule
NGC~6819 & 11.92 & 11.90 & $+0.01$ \\
NGC~2243 & 13.08 & 13.25 & $-0.16$ \\
NGC~7789 & 11.68 & 11.27 & $+0.41$ \\
\bottomrule
\end{tabular}
\caption{Red-clump distance moduli from member giants, against literature
values. Two of the three agree to better than 0.2 mag, which for a membership
list of a few dozen giants assembled from a survey without any prior
membership information is the strongest validation in this chapter: the
members we find are located where the clusters are.}
\label{tab:dm}
\end{table}

\subsection{The case that decided us: NGC~2243 from two surveys}

NGC~2243 is the demonstration that chemical tagging is a two-survey problem.
APOGEE, being a near-infrared survey of the Galactic disc, observes a
cluster's \emph{giants}: bright, plentiful, and excellent for distances and
chemistry. GALAH observes at visible wavelengths and reaches the
\emph{main sequence and turnoff}, which is where age sensitivity lives. Neither
survey alone constrains both.

The sequence we ran is the one a workshop should reproduce:

\begin{enumerate}
\item Take the APOGEE member giants and fit an isochrone. The fit returns an
  age of 1.59 Gyr against a literature value of 1.07 Gyr --- wrong by half ---
  with a broad distance posterior ($\sigma_{dm} = 0.111$).
\item Use the red-clump distance (13.08 mag, agreeing with the literature
  value of 13.25 to within 0.17) as a tight prior, and select GALAH stars
  kinematically around the resulting centroid. At 4.7 kpc the parallax error
  rivals the parallax itself, so a relaxed parallax cut is required --- a
  detail that took a re-run to get right.
\item Fit the combined main-sequence-plus-giants sample: age $0.98$ Gyr
  (residual $0.09$), $\sigma_{dm} = 0.086$.
\end{enumerate}

A factor of six improvement in the age residual, from the same cluster and the
same survey data, by adding the stars that are sensitive to age and prior-ing
the distance with a measurement that is not. The age posterior remains broad
($\sigma_{\log a} \approx 0.97$), which is the degeneracy of \S 15.1 making
itself felt --- but the mode is right, and the distance is a fifth as
uncertain.

\subsection{The finale: better membership, better parameters}

The clearest scientific payoff of the whole project is a single cluster. Fit
an isochrone to three different membership lists for M~67 --- the literature
values are age $2.82$ Gyr, $dm = 9.54$, $[{\rm Fe/H}] = 0.03$:

\begin{table}[htbp]
\centering
\small
\tabcolsep6pt
\begin{tabular}{@{}lccc@{}}
\toprule
membership & $n$ & age (Gyr) & residual \\
\midrule
stage 1 (kinematic candidates, contaminated) & 259 & 1.89 & 0.93 \\
kinematic truth                              & 219 & 2.45 & 0.37 \\
two-stage (spectral rejection)               & 237 & \textbf{2.65} & \textbf{0.17} \\
\bottomrule
\end{tabular}
\caption{Isochrone ages for M~67 from three membership lists, against the
literature age of 2.82 Gyr. The two-stage pipeline --- kinematic candidates,
then rejection of spectral outliers --- recovers the age better than the raw
kinematic membership does.}
\label{tab:m67}
\end{table}

Read the middle row carefully, because it is the surprising one. The kinematic
membership is the \emph{truth} for every benchmark in this workbook, and it
still gets the age wrong by 0.37 Gyr. The two-stage list, which throws away 22
stars that kinematics accepted, gets it wrong by 0.17. The contaminants that
the spectral stage removed were not random: they were young field stars with
the right position and velocity, dragging the fitted age down.

Two honest qualifications. This is one cluster, so it is an existence proof
that membership quality propagates into parameters, not a calibration of the
size of the effect. And the effect is measured in the same direction as the
known bias of contamination on a young cluster's age --- which is reassuring,
and also means the mechanism is understood rather than merely observed.

\subsection{What this does and does not establish}

The Galactic-archaeology promise of chemical tagging is large: if the chemical
fingerprint of a birth cluster survives dispersal into the field, then the
field's stars can be assigned to birth sites, and the disc's star-formation
history can be reconstructed from the field alone, without the kinematic
luck that today's analyses depend on. This workbook measures where that
programme stands at our precision, and the honest summary is:

\begin{itemize}
\item \textbf{Established here.} Spectra contain cluster information that
  element-ratio summaries do not fully capture (\S\ref{sec:spectral}); that
  information is enough to improve membership lists well beyond kinematic
  selection in the specific case we tested (\textbf{Table~\ref{tab:m67}}), and
  it improves a physical parameter as a result; and distances from member
  giants agree with the literature for the clusters where we can measure them.
\item \textbf{Not established.} That the improvement generalises. One cluster,
  one pipeline, and a sample with duplicate rows and a dominant globular
  cluster. A referee will ask for the same experiment on the 25-cluster set
  with five independent membership definitions, and they will be right to.
\item \textbf{Not attempted here.} Recovering a dispersed cluster from the
  field with no kinematic prior. \textbf{Table~\ref{tab:field}} is our best
  attempt and it reaches precision $\approx 0.62$ at recall $0.42$ --- good
  enough to sharpen a membership list, not good enough to find a cluster from
  scratch. Strong chemical tagging remains an open problem in the sense that
  \citet{Casamiquela:21} define it: their best-case sample, with differential
  abundances and no field stars, recovers 9 of 31 clusters at the 40\,\%
  threshold and one at 70\,\%, and their clustering step recovers less than that
  on our stars (\S\ref{sec:literature}).
\end{itemize}

\begin{issues}[EXERCISES]
\begin{enumerate}
\item Reproduce the two-survey fit for one of the other clusters in the
  combined sample. Report age, distance and their posterior widths, with and
  without the red-clump prior. How much of the age improvement comes from the
  main sequence and how much from the prior?
\item The age--metallicity degeneracy means the posterior stays broad even
  with perfect membership. Choose a cluster in the sample and quantify the
  degeneracy directly: plot the posterior in the age--metallicity plane and
  measure its correlation coefficient. What extra observable would break it?
\item Contamination bias has a sign. Construct a synthetic cluster with 10%
  young field stars added, fit it with the same machinery, and compare the
  recovered age with the truth. Repeat with old contaminants. Which is the
  larger effect, and why?
\item \textbf{Table~\ref{tab:m67}} removes 22 stars from the kinematic list
  and improves the age. Inspect those 22 stars: what are their abundances,
  their kinematics and their position relative to M~67? Write the three-sentence
  case for and against calling them contaminants.
\end{enumerate}
\end{issues}
